\section{Реальне застосування: обмежене дослідження цифрових слідів}
\label{sec:application}
Застосування до Петропавлівської Борщагівки показує, як програмне забезпечення працює з інформаційно насиченим локальним конфліктом, не надаючи остаточної реконструкції цього конфлікту. Інцидент і подальші повідомлення охоплюють 2--9 листопада 2026 року; окремі збори мають вужчі вікна. Змістові твердження розглянуто в окремому рукописі кейсу \citep{pylypenko2026case}. Тут застосування організовано за аналітичними передачами: що надійшло до кожного компонента, що він створив і що лишилося нерозв'язаним. Це демонстраційне навантаження, а не новий незалежний бенчмарк інцидентів.

\subsection{Від цільової панелі до зафіксованої рамки аналізу}
Історична панель Telegram має 241 моніторований ідентифікатор каналу та 23\,350 зібраних дописів, із яких 1\,408 із 159 каналів обрано як релевантні. Це цільова панель, а не перепис. Поля тексту й лічильника переглядів у релевантних записах цієї рамки не мають пропусків. Сума їхніх лічильників на момент знімка становить 115\,931\,014; це не розмір унікальної аудиторії й не міра переконання. Вік збору, повторні покази, переслані матеріали й недоступні користувачі унеможливлюють таке тлумачення. Наведення великої суми без правила підрахунку перебільшило б те, що вимірює система.

Пізніша об'єднана рамка містить 23\,386 дописів, 1\,438 релевантних записів і 158 релевантних каналів. Вона має інший SHA-256, ніж історичний файл. Обидві рамки корисні для різних квитанцій аналізу, але їхні знаменники не можна взаємно підставляти. Публікаційний аудит відновив це розрізнення, прочитавши фактичний JSONL і порівнявши його дайджест із входами експериментів. Він також виявив звіт про лінії походження, який приписав точний граф із 369 ребер об'єднаному корпусу всупереч історичним вхідним хешам експерименту. Історичні документи збережено, а атрибуцію виправлено в матриці доказів нової статті. Це демонструє як корисність машинозчитуваного походження, так і потребу перевіряти правильність його тверджень.

\subsection{Від текстових кандидатів до відтворюваних інформаційних сигналів}
Із 1\,408 релевантних записів 1\,384 мають очищений текст довжиною щонайменше 80 символів. Вичерпне міжканальне порівняння Jaccard дає 369 ребер між 357 дописами й 113 каналами. Застарілий результат MinHash містить 367 із цих ребер. Два пропущені ребра мають лаги понад тридцять хвилин, тому їхнє включення змінює знаменник графа, не змінюючи чисельників коротких лагів. Цей результат є конкретним системним уроком: наближений пошук можна оцінити щодо точної обмеженої множини ще до тлумачення статистичних висновків.

На точному графі за порогу Jaccard 0,60 у межах десяти хвилин трапляються 134 пари, у межах тридцяти --- 198. Нульова модель із двогодинним зсувом дає середнє 42,868 пари за десять хвилин і центральний діапазон вибірок 32--55, тоді як 24-годинна модель --- середнє 7,131 і діапазон 2--13. Різниця між цими фонами показує, як умова спільної активності змінює порівняння. Обидві фіксують відхилення в цьому наборі даних; жодна не визначає, хто контролював канали чи чи було копіювання обманним. Свіже відтворення префіксів і арифметичні перевірки оцінюють обчислення, а відсутність підібраних органічних хвиль обмежує змістову специфічність.

Мережа видимих пересилань менша й вимірює інше відношення: 69 кінцевих каналів, 60 спрямованих ребер між парами каналів і 65 спостережених пересилань. Ізольовані канали до цього графа, індукованого ребрами, не входять. Текстова схожість не підміняє явної атрибуції пересилання, а центральність не використовується як доказ організаційного керування. П'ятнадцять шестигодинних наративних інтервалів містять два порожні та останній неповний. Короткий нестаціонарний ряд зі спільною подією та слабка валідація міток не підтримують причинних тверджень про часові ряди. Тому реалізація не надає тлумачення причинності за Ґрейнджером у дослідженому мережевому аналізі.

\subsection{Від обраних обговорень до меж вивчення аудиторії}
Навантаження коментарів містить 39\,598 записів в 533 обраних гілках із 39\,459 доступними ідентифікаторами авторів і 12\,797 псевдонімними авторами. 139 записів без ідентифікатора автора не можуть так само ввійти до аналізу з відомим автором. Набір гілок обрано частково за активністю обговорення. Повнота посторінкового проходу в цих гілках не робить їх репрезентативними щодо всіх коментарів Telegram, мовчазних читачів чи української громадськості. Один автор може давати багато спостережень, тому оцінки зі зважуванням за коментарями й за авторами відповідають на різні запитання.

Можливість отримання реакцій перевірено окремо на десяти коментарях із трьох гілок. Це встановлює, що поточний знімок MTProto може надати подробиці реакцій для цих доступних повідомлень, але не дає виміряного розподілу емодзі в усіх 533 історичних гілках. Контракти аудиторії Social можуть зберігати такі знімки й надані мітки з урахуванням цілі, але наявні докази не виправдовують автоматично згенерованих відсотків гніву, підтримки чи токсичності. Застосування ілюструє важливе інженерне обмеження: збір великого корпусу обговорень не дорівнює наявності валідованої семантичної аналітики для нього.

\subsection{Від документів установ до обмежених станів доказів}
Спеціалізований збирач тверджень надає 25 обраних суджень і 25 обраних уривків. Поточний результат містить 19 підтриманих у межах обсягу та шість із недостатніми доказами, спростованих у межах обсягу немає. Ці стани можуть документувати, що установа зробила заяву, залишаючи нерозв'язаним фізичне звинувачення, про яке йдеться. Зокрема, офіційні заперечення зберігаються як офіційні позиції, а не як незалежний доказ того, що фізичної події не було. Корені джерел і типи обсягу лишаються наданими інтерпретаційними метаданими; їх наявність у схемі не перевіряє незалежності чи правильності обсягу.

Застосування також показує межу лексичних правил відношень: уривки про травми приймалися для судження про самовільне залишення частини, бо збігалося число. Це істотний збій попри дійсні посилання й числовий збіг. Він не дозволяє вважати 19 підтриманих результатів 19 незалежно встановленими істинами. Післяаудитна ревізія з розділу~\ref{sec:rule-v2} відкликає ці прийняття, а також утримує твердження про судову заставу, уривок якого не називає особу, що лишає 18 підтриманих і сім із недостатніми доказами. Ці 18 так само не є незалежно встановленими істинами. Загальний запускач обирає консервативніший шлях: пошук кандидатів без прийнятих оцінок. Це різні шляхи коду, і стаття не переносить інвентар кураторського збирача на загальний процес.

\subsection{Від байтів медіа до кандидатних спостережень}
Історичні архівні квитанції перелічують 775 рядків дописів-джерел із перевіреним завантаженням, що відповідають 607 унікальним активам SHA-256. Суворий технічний індекс фіксує всі 607 як декодовані, загальним обсягом 9\,472\,884\,834 байти. Квитанції обробки описують 8\,679 кадрів, наявність аудіо у 602 активах і його відсутність у п'яти та приблизно 12,04 сумарної години тривалості контейнерів. Тривалість контейнера не є тривалістю мовлення; ці історичні вимірювання для цієї публікації не повторювали.

Масовий ASR створив 10\,872 кандидатні сегменти в 558 повних текстових результатах, 44 результати без сегментів і п'ять активів без аудіо. 18 обраних еталонних кліпів дають 185 помилок на 415 еталонних слів. Ці інвентарі показують масштаб технічної обробки й потенційну помилку, а не автентичність, особу мовця чи акустичну перевірку. Кандидат візуального збігу може спрямовувати огляд аналітиком, але не може встановити, що кліп зображує прив'язане до нього локальне звинувачення. Застосування завершується структурованими невідомими скрізь, де автентичність джерела, місце, фізичні події чи семантичне випливання лишаються невалідованими. Його практична цінність --- простежуваний аналітичний робочий простір, а не остаточне вирішення конфлікту.
